\documentclass[%
 superscriptaddress,
 amsmath,amssymb,
 prx,
]{revtex4-2}
\usepackage{graphicx} % Required for inserting images
\usepackage{comment}
\usepackage{dsfont}
\usepackage{color}
\usepackage{bm} % bold math with {\bm ...text...} command
\usepackage{hyperref}

\usepackage[utf8]{inputenc}
\usepackage{amsmath}
\usepackage{amssymb}
\usepackage[a4paper, total={6.6in, 8.8in}]{geometry}
\usepackage{graphicx}
\usepackage{hyperref}
\usepackage{float}
\usepackage{braket}
\usepackage{adjustbox}
\usepackage{color}
\usepackage{framed} 
\usepackage{url}
\usepackage[normalem]{ulem}

\usepackage{color}
\definecolor{eggplant}{RGB}{180,33,147}
\newcommand{\mf}[1]{{\color{blue}{[MF: #1]}}}
\newcommand{\jt}[1]{{\color{red}{[JT: #1]}}}
\newcommand{\ds}[1]{{\color{green}{[DS: #1]}}}
\newcommand{\msc}[1]{{\color{eggplant}{[MS: #1]}}}
\newcommand{\ms}[1]{{\color{eggplant}{#1}}}

\begin{abstract}
We report an accurate, memory and time efficient classical simulation of a 127-qubit kicked Ising quantum system on the heavy-hexagon lattice. A simulation of this system on a quantum processor was recently performed using noise mitigation techniques to enhance accuracy (Nature volume 618, p.~500–505 (2023)). Here we show that, by adopting a tensor network approach that reflects the qubit connectivity of the device, we can perform a classical simulation that is significantly more accurate than the results obtained from the quantum device in the verifiable regime and comparable to the quantum simulation results for larger depths. The tensor network approach used will likely have broader applications for simulating the dynamics of quantum systems with tree-like correlations.
\end{abstract}


\begin{document}

\title{Efficient tensor network simulation of IBM's Eagle kicked Ising experiment}


\author{Joseph Tindall}
\affiliation{Center for Computational Quantum Physics, Flatiron Institute, New York, New York 10010, USA}

\author{Matthew Fishman}
\affiliation{Center for Computational Quantum Physics, Flatiron Institute, New York, New York 10010, USA}

\author{E. Miles Stoudenmire}
\affiliation{Center for Computational Quantum Physics, Flatiron Institute, New York, New York 10010, USA}

\author{Dries Sels}
\affiliation{Center for Computational Quantum Physics, Flatiron Institute, New York, New York 10010, USA}
\affiliation{Center for Quantum Phenomena, Department of Physics, New York University, 726 Broadway, New York, NY, 10003, USA}



\maketitle


\section*{Main}
\par \emph{Introduction.} Quantum computers are fundamentally noisy in nature and incur errors with each operation. To remedy noise and make the computer run as desired, one option could be quantum error correcting codes. But in spite of significant advances in quantum technologies over the last decade, this type of error correction is currently out of reach. As such, huge efforts have been devoted to finding out whether current noisy quantum technologies could provide a practical advantage over classical computers without error correction. The emphasis here is on practicality, since expectation values in shallow quantum circuits can formally be computed in polynomial time on a classical computer~\cite{bravyi21,wild23} and so can the output distribution of a noisy random quantum circuit~\cite{aharonov23}. The latter refutes the claim of violation of the extended Church-Turing thesis in Ref.~\cite{arute19}, and makes it exceedingly unlikely any such violation can be achieved through noisy quantum computation. While these works put formal constraints on the power of noisy quantum devices, they also highlight that it may be difficult to classically simulate them in practice because the polynomials have rather unfavorable scaling when the noise is sufficiently small. At the same time, error mitigation techniques have been proposed to further improve the accuracy of noisy quantum devices by classical post processing~\cite{temme17,vandenberg23}. With these techniques one can construct a better classical (biased) estimator for the desired noiseless result by paying the overhead of having to run several noisy quantum circuits.

The power of these techniques was recently demonstrated in the context of dynamical quantum simulations of a two-dimensional (2D) transverse-field Ising model on a $127$-qubit heavy-hexagon lattice~\cite{kim2023} (see Fig.~\ref{Fig:heavyhexdiagram}). This lattice corresponds to the connectivity of the quantum processor (IBM Eagle) used in that study. The dynamics were implemented in Trotterized fashion, resulting in so-called `kicked Ising' dynamics given by the unitary: 
\begin{equation}
    U(\theta_{h}) = \left(\prod_{\langle v,v' \rangle}\exp \left({\rm i} \frac{\pi}{4}Z_{v}Z_{v'} \right)\right)\,\left(\prod_{v}\exp \left(-{\rm i} \frac{\theta_{h}}{2}X_{v}\right)\right),
    \label{Eq:Propagator}
\end{equation}
where $Z$ and $X$ denote Pauli operators and $\langle v,v' \rangle$ indicates that $v$ and $v'$ are neighbours on the heavy-hex lattice. Expectation values of a number of observables were extracted after a few Floquet cycles using zero-noise extrapolation techniques~\cite{temme17}. 
The experimental results were found to be much more accurate than several classical tensor network approaches
applied to the same problem, even when the tensor network methods utilized significant computational resources. Here we show that, by adopting a different tensor network approach that respects the qubit connectivity of the heavy-hex lattice  (see Fig.~\ref{Fig:heavyhexdiagram}), that we can perform classical simulations of this problem with modest computational resources and to a much higher degree of accuracy than the previously reported classical results or the error mitigated quantum computer in the verifiable regime. For larger depth circuits, we achieve an accuracy comparable to the quantum processor results. The tensor network method we use exploits the locally tree-like nature of the lattice, and demonstrates the feasibility of using tensor networks for studying the dynamics of quantum systems on tree-like graphs. Before describing the method, let us present the results.
\par \emph{Results.} In Fig.~\ref{Fig:comparison} we compare our method to the results presented in Fig.~3 of Ref.~\cite{kim2023}. The availability of exact results in this case, based on brute force light-cone simulation techniques, allows us to directly assess the errors. The tensor network state (TNS) and optimization methods we use simulate the full $127$-qubit system and result in highly accurate expectation values --- the expected value of the average single-site magnetization can be computed to an accuracy of $\sim 10^{-14}$ with a simulation that runs in less than $2$ minutes, see Fig.~\ref{Fig:comparison}(a). In this classically verifiable regime we systematically outperform the quantum processor results reported in Ref.~\cite{kim2023} to orders of magnitude and with very modest resources. Specifically, for every expectation value obtained by the processor and plotted in Fig. 3 of Ref. \cite{kim2023} we obtain a corresponding value to orders of magnitude better accuracy with a simulation that takes less than $7$ minutes to run on a single Intel Skylake CPU and a state that takes up, at most, $0.3$GB of memory. 


\begin{figure}[t!]
    \centering
    \includegraphics[width = \columnwidth]{Figures/HeavyHexTNSZoomed.pdf}
    \caption{Qubit layout for the Eagle quantum processor and the tensor network structure used for our simulations. The structure directly reflects the processor layout. On-site tensors $\Gamma_{v}$ are coloured in blue and possess physical, uncontracted, indices of dimension $2$ (represented by their dangling legs) and virtual indices of dimension $\chi$ (represented by the edges of the network) which are shared with neighboring tensors. Positive, diagonal bond tensors $\Lambda_{e}$ live on the edges $e$ between the site tensors and are coloured in grey.}
    \label{Fig:heavyhexdiagram}
\end{figure}

\begin{figure}[t!]
    \centering
    \includegraphics[width = \columnwidth]{Figures/Figure1.pdf}
    \caption{Comparison for classically verifiable systems of our tensor network state using the belief propagation approximation (BP-TNS) approach to simulating the dynamics of the 2D transverse field Ising model versus the Eagle quantum processor and alternative tensor network methods. Expectation values with respect to the state $\vert \psi(\theta_{h}, 5) \rangle$ (i.e following $5$ Trotter steps of the dynamics of the model --- see Eqs. (\ref{Eq:Propagator}) and (\ref{Eq:State})) are plotted, alongside exact results determined from light-cone simulations. \textbf{a)} Average magnetization. \textbf{b)} Weight-10 observable. \textbf{c)} Weight-17 observable. The bottom plots show errors defined as the absolute difference between the simulation result and the exact result. For some data points the error from our BP-TNS simulation is too small to fit on the scale of the plot and so these points are not marked. Circled, annotated points denote, for a given $\theta_{h}$, the memory required to store the final state and the walltime associated with performing the simulation and calculating the relevant observable on a single Intel Skylake CPU.}
    \label{Fig:comparison}
\end{figure}

\par In Fig.~\ref{Fig:F2} we compare our method to the results presented in Fig.~4 of Ref.~\cite{kim2023}, showing properties after 
larger numbers of Trotter evolution steps. 
In this regime, the increased depth of the circuit means exact data is no longer available for benchmarking. For both the $n = 6$ and $n = 20$ Trotter step simulations we capture the known exact values at the Clifford points $\theta_{h} = 0$ and $\theta_{h} = \pi/2$ --- with the weight-$17$ stabilizer demonstrating the anticipated saturation to unity that is not captured by the matrix product state (MPS) approach. 
Moreover, in Fig.~\ref{Fig:F2}c we show the convergence of both observables with the bond dimension of our TNS at the specified values of $\theta_{h}$, performing an extrapolation to infinite bond dimension which demonstrates the closeness of our finite-$\chi$ results to the limit $\chi \rightarrow \infty$ at the chosen $\theta_{h}$ values. These are generic points where the expectation value has a finite, non-trivial value.
An extrapolation in $\chi$ in the regime $\pi/4 \leq \theta_{h} \leq 7\pi/16$ in Fig. \ref{Fig:F2}b would likely lead to the most significant changes with bond dimension and our results in this regime could differ slightly from the $\chi \rightarrow \infty$ limit.
While expectation values in our simulation are calculated approximately, we have seen in Fig. \ref{Fig:comparison} that this approximation is extremely accurate, at least for shorter circuit depths. For larger circuit depths, such as in Fig. \ref{Fig:F2}b,  we are making an uncontrolled approximation and so even in the limit $\chi \rightarrow \infty$ the results will differ from the exact answer. The accuracy of this could be quantified with more controlled methods that go beyond rank-one approximations of the environments, as we discuss in the conclusion.

\par The accuracy of our results here can be understood from the locally tree-like nature of the heavy-hex lattice, with the smallest loop having a path length of $12$. Our chosen methods for performing the time evolution and taking expectation values (see Methods) are optimal if the network is a tree. The heavy-hex lattice is close enough to a tree for such methods to still be highly effective. 

\par \emph{State ansatz.} Our ansatz for the wave function of the system is a tensor network which directly reflects the `heavy-hex' qubit connectivity of the processor (see Fig.~\ref{Fig:heavyhexdiagram}). The physical properties of a tensor network are invariant under a gauge symmetry corresponding to the insertion of any invertible matrix $G_{e}$ and its inverse $G_{e}^{-1}$ on any contracted, internal bond/edge $e$ of the network.
We take advantage of this symmetry to keep the tensor network in the \emph{Vidal gauge} \cite{Vidal2003, Jiang2008}. This gauge corresponds to the choice of positive, diagonal bond tensors $\Lambda_{e}$ residing on the edges of the network and the on-site tensors $\Gamma_{v}$ of the network obeying certain isometric properties (see Eqs. (\ref{Eq:IsometryPropertyDefinition}) and (\ref{Eq:IsometryDefinition})). These isometric properties are important for maintaining accuracy during the evolution of the network. We use $\vert \psi(\theta_{h}, n) \rangle$ to denote the TNS of the system after $n \geq 1$ applications of $U(\theta_{h})$, i.e. 
\begin{equation}
    \vert \psi(\theta_{h}, n) \rangle = \left(\prod_{i=1}^{n}U(\theta_{h})\right)\vert \psi(0) \rangle = U^{n}(\theta_{h})\vert \psi(0) \rangle,
    \label{Eq:State}
\end{equation}
where $\vert \psi(0) \rangle$ is the initial state of the system. We use the same initial state as in Ref. \cite{kim2023}: $\vert \psi(0) \rangle = \vert\! \uparrow \uparrow \hdots \uparrow \rangle$.
The single-site $X$ rotations in $U(\theta_{h})$ can be applied to $\vert \psi(\theta_{h}, n) \rangle$ exactly and the two-site gates are applied approximately using the simple update \cite{Jiang2008} procedure (see Methods section), which involves truncating the internal indices of the TNS to keep them less than or equal to a prescribed maximum bond dimension $\chi$. Following a single Trotter step, the tensor network is re-gauged using belief propagation, a well-established statistical inference algorithm which can  be formulated for tensor networks \cite{TensorNetworkBeliefPropagation2}. Belief propagation allows us to rapidly determine the necessary transformation matrices by performing \emph{message passing} over the network representing $\langle \psi(\theta_{h}, n) \vert \psi(\theta_{h}, n) \rangle$. The use of belief propagation as a method to efficiently `re-gauge' a tensor network was recently formulated by some of the current authors in Ref. \cite{TindallGauging2023} and is closely related to other known gauging methods  \cite{Ran2012,Phien2015,Jahromi2019,SimpleUpdateGauging3}. We emphasize that the results presented here could have been achieved with those known gauging methods.

\par \emph{Mean values.} We measure single-site expectation values of the gauged state $\vert \psi(\theta_{h}, n) \rangle$ using a rank-one approximation for the environments of local regions of the network. Such a method is only guaranteed to be exact in the limit of a tree network. Nonetheless, as demonstrated by the results in Fig.~\ref{Fig:comparison}, the large-loop structure of the 
heavy-hexagon lattice means that the network is locally tree-like and therefore amenable to such an approximation. In Ref.~\cite{kim2023} specific higher-weight observables were also measured to highlight that non-trivial results can be obtained in the ``regime of strong entanglement'' --- such observables can be difficult to accurately measure for loopy tensor networks. 
We can, however, exploit the Clifford properties of the circuit at $\theta_{h} = \frac{\pi}{2}$ to transform the problem of measuring higher-weight observables after $n$ Trotter steps into one of measuring a single-site observable after evolution by $n$ Trotter steps with the propagator $U(\theta_{h})$ and $n$ Trotter steps with the propagator $\left(U(\frac{\pi}{2})\right)^{\dagger}$. 
Specifically, the operator $U^{n}\left(\frac{\pi}{2}\right)Z_{i}\left(U^{n}\left(\frac{\pi}{2}\right)\right)^{\dagger}$ is always a single Pauli string and its expectation with respect to $\vert \psi(\theta_{h}, n) \rangle$ can be obtained by further evolving $\vert \psi(\theta_{h}, n) \rangle$ with $\left(U(\frac{\pi}{2})\right)^{\dagger}$ and then measuring $Z_{i}$. For instance,
\begin{equation}
    X_{13, 29, 31}Y_{9,30}Z_{8,12,17,28,32} =  U^{5}\left(\frac{\pi}{2}\right)Z_{13}\left(U^{5}\left(\frac{\pi}{2}\right)\right)^{\dagger},
\end{equation} 
meaning this observable can be calculated, with respect to $\vert \psi(\theta_h,5) \rangle$, by measuring $\langle Z_{13} \rangle$ with respect to the state $\left( U^{5}\left(\frac{\pi}{2}\right)\right)^{\dagger}\vert \psi(\theta_h,5) \rangle$. We can reach this state straightforwardly by performing a $10$ ($5$ + $5$) Trotter step tensor network simulation. We use this `extended time-evolution' method to measure all higher-weight observables presented in Ref. \cite{kim2023} and note that this procedure is entirely generic: any long Pauli string can be generated out of a single Pauli by application of a Clifford circuit. We would also like to emphasize that these higher-weight observables could be measured with other tensor network methods. For instance, the planar nature of our heavy-hex TNS means a generalized boundary MPS \cite{Verstraete2004} method such as `sweep line contraction' \cite{chubb2021general} can be directly applied to the state $|\psi(\theta_h,n) \rangle$ to measure a desired higher-weight observable.

\begin{figure}[t!]
    \centering
    \includegraphics[width = \columnwidth]{Figures/Figure2.pdf}
    \caption{Comparison for non-classically-verifiable systems of our tensor network state using the belief propagation approximation (BP-TNS) approach to simulating the dynamics of the 2D transverse field Ising model versus the Eagle quantum processor and alternative tensor network methods. Expectation values calculated following a number of Trotter steps of the dynamics of the model --- see Eq. (\ref{Eq:Propagator}) --- are plotted. \textbf{a)} Weight-17 stabilizer after $6$ steps of evolution. \textbf{b)} Weight-1 observable after $20$ steps. \textbf{c)} Top and bottom correspond to scaling, with inverse bond dimension of the heavy-hex TNS ansatz, of the observables in  \textbf{a)} and \textbf{b)} respectively at select values of $\theta_{h}$. The matching data points are circled. Red-dashed lines represent a least-squares fit of the form $A + B \slash \chi$ taken on the data for $\chi \geq 150$ --- with $A$, the predicted value of the observable in the limit $\chi \rightarrow \infty$, being $0.2736$ and $0.4788$ (to $4$ significant figures) for the top and bottom plots respectively. Even in the limit $\chi \rightarrow \infty$ there will be some deviation from the exact result due to our rank-one method for taking expectation values (see Methods).}
    \label{Fig:F2}
\end{figure}

\par \emph{Conclusion.} 
We have shown that a 127-qubit simulation of Floquet dynamics of the kicked Ising model on a heavy-hexagon lattice, recently simulated on a quantum processor in Ref.~\cite{kim2023}, can be performed accurately and with minimal computational resources with tensor networks.
Our work  demonstrates the importance of adopting a tensor network ansatz which  reflects the spatial connectivity and entanglement structure of the system. The chosen optimization method for this ansatz, based on simple update gate evolution and belief propagation, is highly effective as it directly exploits the locally tree-like nature of the lattice. The MPS ansatz used for benchmarking in Ref. \cite{kim2023} requires large bond dimensions even for a depth-$5$ circuit due to the Hamiltonian becoming long-range when mapped onto one dimension. Meanwhile, the isoTNS ansatz suffers from the drawback of a significant $\mathcal{O}(\chi^{7})$ scaling with bond dimension and it is a newer method whose best practices and representational power are still under investigation \cite{Zaletel2020, Soejima2020}. Our ansatz directly reflects the heavy-hex geometry and, for the problem at hand, scales as $\mathcal{O}(L \chi^{4})$ for a given Trotter step evolution, where $L$ is the number of qubits and $\chi$ is the bond dimension. Generically, the bond dimension $\chi$ will grow as more Trotter steps are applied.

An important takeaway from our results is that there can be many complementary routes toward classical simulation of quantum many-body systems, especially for those with physical structure such as separation of energy scales or locality. In addition to well-known properties such as low entanglement, shallow circuit depth, or low T-gate count (in the case of nearly Clifford circuits), our work highlights the importance of network structure as another key property that can be exploited. The methods we have adopted here directly take advantage of the locally tree-like properties of the heavy-hexagon lattice. 
For quantum states with less tree-like correlations, including the state on the heavy-hex lattice at larger circuit depths, it is nevertheless possible to obtain more accurate results by going beyond rank-one approximations of environments with methods like boundary MPS contractions of 2D tensor networks \cite{Verstraete2004, Murg2007}, while still having the TNS structure match the processor design. Quantifying the errors coming from rank-one approximations with these more controlled methods is an important area for future study.  
We also emphasize that tensor network methods are not limited to one- and two-dimensional systems, and higher-dimensional, non-planar systems can actually become more mean-field-like, allowing tensor network approaches to again work well~\cite{tindall2022quantum, searle2023exact}. 

\par Our work opens up new directions in which highly flexible and computationally inexpensive tensor network approaches can be used to benchmark new quantum processor designs and can better delineate which many-body quantum systems could become difficult for classical computing techniques. The software we used is part of the forthcoming ITensorNetworks.jl package \cite{ITensorNetworks} being developed at the Flatiron Institute  (see Computing Resources and Software Packages), which enables rapid testing and deployment of tensor network methods on arbitrary graphs. This software is continually being developed with the aim of tackling the type of simulation presented here.  


\section*{Methods}
\label{sec:methods}
\subsection*{The Vidal gauge}
Our tensor network ansatz for the state of the system consists of both local tensors $\Gamma_{v}$ on the sites of the network and bond tensors $\Lambda_{e}$ which live on the edges of the network. The state is in the `Vidal' gauge which is characterised by a set of isometric constraints on certain groupings of local tensors. Specifically, for every edge $e = (v, v_{i})$ connecting vertices $v$ and a neighboring vertex $v_i$, if we group the following tensors
\begin{equation}
    \adjincludegraphics[valign=c]{Diagrams/W.pdf} \quad,
    \label{Eq:IsometryPropertyDefinition}
\end{equation}
they form an isometric tensor obeying
\begin{equation}
    \adjincludegraphics[valign=c]{Diagrams/IsometryPropertyTNS.pdf} \quad ,
    \label{Eq:IsometryDefinition}
\end{equation}
where the RHS denotes the identity matrix and we have used the set $\{v_{1}, v_{2}, \hdots v_{d}\}$ to denote the $d$ neighbors of vertex $v$. Having the TNS obey this property is important for maintaining accuracy when applying two-site gates and when taking expectation values. Below we detail the methods we use for performing these operations. We then detail how we use belief propagation to maintain the Vidal gauge during our simulations.

\subsection*{The simple update procedure}
To apply two-site gates to the tensor network state we adopt the simple update procedure \cite{Jiang2008}. The procedure is depicted diagramatically in Fig. \ref{fig:SimpleUpdate}. For better efficiency, we use the `reduced tensor' variant of simple update \cite{Wang2011,Li2012} (not shown in Fig. \ref{fig:SimpleUpdate}).
\begin{figure}[t!]
    \centering
    \includegraphics[width =\columnwidth]{Diagrams/SimpleUpdate.pdf}
    \caption{Steps to perform a simple update on the edge of a TNS in Vidal gauge. The example pictured is two neighboring sites of degree $2$ and $3$ respectively. \textbf{i) - ii)} The bond tensors and gate are combined into a single composite tensor $\Theta_{v,w}$. \textbf{ii) - iii)} An SVD is performed on $\Theta_{v,w}$ and singular values of the resulting bond tensor $\tilde{\Lambda}_{v,w}$ can be discarded in order to truncate the bond. \textbf{iii) - iv)} Resolutions of identity, using the original bond tensors, are inserted on the exposed edges. \textbf{iv) - v)} The inverse bond tensors are absorbed, resulting in the updated local tensors $\tilde{\Gamma}_{v}$ and $\tilde{\Gamma}_{w}$.}
    \label{fig:SimpleUpdate}
\end{figure}

\subsection*{Measuring expectation values within the Vidal gauge}
In order to measure a single-site observable $\langle O_{v} \rangle$ on site $v$ of our TNS we absorb the neighboring bond tensors onto the on-site tensor $\Gamma_{v}$ and contract the result with its conjugate, inserting the single-site observable $O_{v}$ along the physical index which is being contracted over.
For the example of a site with $3$ neighbours this can be visualized as
\begin{equation}
    \adjincludegraphics[valign=c]{Diagrams/ApproxSingleSiteExpec.pdf} \quad,
    \label{Eq:SingleSIte}
\end{equation}
where we have defined the square of the bond tensors
\begin{equation}
    \adjincludegraphics[valign=c]{Diagrams/SquareBondTensor.pdf} \quad.
    \label{Eq:SquareBondTensors}
\end{equation}
Equation (\ref{Eq:SingleSIte}) approximates $\langle O_{v} \rangle$ due to not accounting for the exact environment surrounding $v$ and instead treating it as a tensor product of environments coming from each of the neighbors of $v$. This is sufficient to get a highly accurate approximation for single-site observables associated with the BP-TNS due to its locally tree-like structure. 
\par In the section below we detail how we use belief propagation to maintain the Vidal gauge during our simulations. This helps maintain accuracy in the simple update procedure and when taking expectation values.


\subsection*{Belief propagation on a tensor network state}
At the heart of our method for maintaining the gauge properties of our tensor network state is belief propagation (BP). BP is a well-established technique for approximating the marginals of the probability distributions of graphical models \cite{BeliefPropagation1} and has recently gained interest in the context of contracting tensor networks \cite{TensorNetworkBeliefPropagation1,TensorNetworkBeliefPropagation2,TensorNetworkBeliefPropagation3,SimpleUpdateGauging3}. 
\par To perform belief propagation on a TNS $\vert \psi \rangle$ in the Vidal gauge we first absorb the square roots of the bond tensors onto the $\Gamma_{v}$ tensors. For instance, for the example of a tensor with $3$ neighbors we define
\begin{equation}
    \adjincludegraphics[valign=c]{Diagrams/SymmetricGaugeTensor.pdf} \quad.
    \label{Eq:SymmetricTensor}
\end{equation}
We then form the closed network which respresents the square norm $\langle \psi \vert \psi \rangle$ of $\vert \psi \rangle$. This `norm' network consists of the on-site tensors $\mathcal{T}_{v} = \sum_{s_v} (T^{s_v}_{v})(T^{s_v}_{v})^*$, where the summation is over the external indices $s_{v}$ of the tensors $T_{v}$ and $\left( T_{v} \right)^{*}$. 
\par Following this, we define a series of `message tensors' over the the norm network, with $M_{v, v_{i}}$ denoting the message tensor directed along the edge from $\mathcal{T}_{v}$ to its neighbor $\mathcal{T}_{v_{i}}$. The indices of $M_{v, v_{i}}$ match the indices shared by tensors $\mathcal{T}_{v}$ and $\mathcal{T}_{v_{i}}$. Here, the direction of the edge is important and, generally, $M_{v,v_{i}} \neq  M_{v_{i},v}$.
A set of self-consistent equations is defined for the message tensors:
\begin{equation}
    M_{v, v_{i}} = \left(\prod_{j \in \{1, 2, ..., d \} \setminus \{i \}}M_{v_{j}, v}\right)\mathcal{T}_{v}
    \label{Eq:MessageTensorSelfConsistency}
\end{equation}
where the product indexes all $d$ neighbours of $v$ excluding $v_{i}$ and multiplication of two tensors implies a contraction over any common indices they share.
This equation can be expressed diagrammatically as
\begin{equation}
    \adjincludegraphics[valign=c]{Diagrams/MessageTensorUpdate.pdf} \quad.
    \label{Eq:BPEquation}
\end{equation}
Initializing the message tensors, one can iterate these equations in an attempt to converge them. The converged messages then form a rank-one approximation of the exact environment for a given tensor $\mathcal{T}_{v}$, and can be used to approximate expectation values of the TNS. The more tree-like the network, the better the approximation. Importantly, these message tensors can then be used to define a gauge transformation which, when applied to a TNS, guarantee satisfaction of Eqs. (\ref{Eq:IsometryPropertyDefinition}) and (\ref{Eq:IsometryDefinition}) along all edges of the network. 
A detailed description of this method, including extensive benchmarking and discussion on its relation to other gauging methods, is given in Ref. \cite{TindallGauging2023}. 
\par In our simulations here we perform belief propagation gauging after every Trotter step. This is to maintain accuracy in these procedures. We should note that for shorter depth circuits (such as those simulated in Fig. \ref{Fig:comparison}) gauging after every Trotter step does not have a significant affect on accuracy. The circuit here is not deep enough to significantly alter the isometric condition in Eq. (\ref{Eq:IsometryDefinition}) and affect the simple update procedure. For longer depth circuits, such as the one run in Fig. \ref{Fig:F2}b, gauging every Trotter step is important and not performing gauging during the simulation leads to a significant loss in accuracy. 
\par We also always gauge the state before taking expectation values, which is equivalent to computing the expectation values with BP. Even for the shorter circuits in Fig. \ref{Fig:comparison}, not performing gauging prior to taking an expectation value can noticeably affect the accuracy of the result --- unless the bond dimension used is large enough for the simulation to be exact and so the simple update procedure preserves the gauge properties of the TNS.

\subsection*{Computing Resources and Software Packages}
The code used to produce the numerical results in this paper was written using the \textbf{ITensorNetworks.jl} package \cite{ITensorNetworks} --- a general purpose and publicly available Julia \cite{bezanson2017julia} package for manipulating (gauging, contracting, partitioning, evolving, etc.) tensor network states of arbitrary geometry. It is built on top of \textbf{ITensors.jl} \cite{itensor-r0.3}, which provides the basic tensor operations. 
Code is available in the current version of ITensorNetworks.jl for performing belief propagation, gauging, and the simple update procedure on arbitrary tensor network states. All code was run using the Flatiron Institute's Rusty computing cluster. Timings and memory usage are based on code being run on $6$ cores of an Intel Skylake node.
\par The tensor network diagrams in this paper were produced using the publicly available package \textbf{GraphTikz.jl} \cite{GraphTikz}, a general-purpose Julia package for visualizing graphs, including tensor networks.

\section*{Acknowledgements}
J.T., M.F., M.S., and D.S. are grateful for ongoing support through the Flatiron Institute, a division of the Simons Foundation. D.S. was supported by AFOSR: Grant FA9550-21-1-0236. We would like to acknowledge Michael Zalatel, Sajant Anand and Bo Xiao for insightful discussions. We also would like to acknowledge Ariel Stolbun, Andy Millis, and Antoine Georges for helpful comments on the draft.

%apsrev4-2.bst 2019-01-14 (MD) hand-edited version of apsrev4-1.bst
%Control: key (0)
%Control: author (8) initials jnrlst
%Control: editor formatted (1) identically to author
%Control: production of article title (0) allowed
%Control: page (0) single
%Control: year (1) truncated
%Control: production of eprint (0) enabled
\begin{thebibliography}{31}%
\makeatletter
\providecommand \@ifxundefined [1]{%
 \@ifx{#1\undefined}
}%
\providecommand \@ifnum [1]{%
 \ifnum #1\expandafter \@firstoftwo
 \else \expandafter \@secondoftwo
 \fi
}%
\providecommand \@ifx [1]{%
 \ifx #1\expandafter \@firstoftwo
 \else \expandafter \@secondoftwo
 \fi
}%
\providecommand \natexlab [1]{#1}%
\providecommand \enquote  [1]{``#1''}%
\providecommand \bibnamefont  [1]{#1}%
\providecommand \bibfnamefont [1]{#1}%
\providecommand \citenamefont [1]{#1}%
\providecommand \href@noop [0]{\@secondoftwo}%
\providecommand \href [0]{\begingroup \@sanitize@url \@href}%
\providecommand \@href[1]{\@@startlink{#1}\@@href}%
\providecommand \@@href[1]{\endgroup#1\@@endlink}%
\providecommand \@sanitize@url [0]{\catcode `\\12\catcode `\$12\catcode
  `\&12\catcode `\#12\catcode `\^12\catcode `\_12\catcode `\%12\relax}%
\providecommand \@@startlink[1]{}%
\providecommand \@@endlink[0]{}%
\providecommand \url  [0]{\begingroup\@sanitize@url \@url }%
\providecommand \@url [1]{\endgroup\@href {#1}{\urlprefix }}%
\providecommand \urlprefix  [0]{URL }%
\providecommand \Eprint [0]{\href }%
\providecommand \doibase [0]{https://doi.org/}%
\providecommand \selectlanguage [0]{\@gobble}%
\providecommand \bibinfo  [0]{\@secondoftwo}%
\providecommand \bibfield  [0]{\@secondoftwo}%
\providecommand \translation [1]{[#1]}%
\providecommand \BibitemOpen [0]{}%
\providecommand \bibitemStop [0]{}%
\providecommand \bibitemNoStop [0]{.\EOS\space}%
\providecommand \EOS [0]{\spacefactor3000\relax}%
\providecommand \BibitemShut  [1]{\csname bibitem#1\endcsname}%
\let\auto@bib@innerbib\@empty
%</preamble>
\bibitem [{\citenamefont {Bravyi}\ \emph {et~al.}(2021)\citenamefont {Bravyi},
  \citenamefont {Gosset},\ and\ \citenamefont {Movassagh}}]{bravyi21}%
  \BibitemOpen
  \bibfield  {author} {\bibinfo {author} {\bibfnamefont {S.}~\bibnamefont
  {Bravyi}}, \bibinfo {author} {\bibfnamefont {D.}~\bibnamefont {Gosset}},\
  and\ \bibinfo {author} {\bibfnamefont {R.}~\bibnamefont {Movassagh}},\
  }\bibfield  {title} {\bibinfo {title} {Classical algorithms for quantum mean
  values},\ }\href {https://doi.org/10.1038/s41567-020-01109-8} {\bibfield
  {journal} {\bibinfo  {journal} {Nature Physics}\ }\textbf {\bibinfo {volume}
  {17}},\ \bibinfo {pages} {337} (\bibinfo {year} {2021})}\BibitemShut
  {NoStop}%
\bibitem [{\citenamefont {Wild}\ and\ \citenamefont {Alhambra}(2023)}]{wild23}%
  \BibitemOpen
  \bibfield  {author} {\bibinfo {author} {\bibfnamefont {D.~S.}\ \bibnamefont
  {Wild}}\ and\ \bibinfo {author} {\bibfnamefont {A.~M.}\ \bibnamefont
  {Alhambra}},\ }\bibfield  {title} {\bibinfo {title} {Classical simulation of
  short-time quantum dynamics},\ }\href
  {https://doi.org/10.1103/PRXQuantum.4.020340} {\bibfield  {journal} {\bibinfo
   {journal} {PRX Quantum}\ }\textbf {\bibinfo {volume} {4}},\ \bibinfo {pages}
  {020340} (\bibinfo {year} {2023})}\BibitemShut {NoStop}%
\bibitem [{\citenamefont {Aharonov}\ \emph {et~al.}(2023)\citenamefont
  {Aharonov}, \citenamefont {Gao}, \citenamefont {Landau}, \citenamefont
  {Liu},\ and\ \citenamefont {Vazirani}}]{aharonov23}%
  \BibitemOpen
  \bibfield  {author} {\bibinfo {author} {\bibfnamefont {D.}~\bibnamefont
  {Aharonov}}, \bibinfo {author} {\bibfnamefont {X.}~\bibnamefont {Gao}},
  \bibinfo {author} {\bibfnamefont {Z.}~\bibnamefont {Landau}}, \bibinfo
  {author} {\bibfnamefont {Y.}~\bibnamefont {Liu}},\ and\ \bibinfo {author}
  {\bibfnamefont {U.}~\bibnamefont {Vazirani}},\ }\bibfield  {title} {\bibinfo
  {title} {A polynomial-time classical algorithm for noisy random circuit
  sampling},\ }in\ \href {https://doi.org/10.1145/3564246.3585234} {\emph
  {\bibinfo {booktitle} {Proceedings of the 55th Annual ACM Symposium on Theory
  of Computing}}},\ \bibinfo {series and number} {STOC 2023}\ (\bibinfo
  {publisher} {Association for Computing Machinery},\ \bibinfo {address} {New
  York, NY, USA},\ \bibinfo {year} {2023})\ p.\ \bibinfo {pages}
  {945–957}\BibitemShut {NoStop}%
\bibitem [{\citenamefont {Aurate}\ \emph {et~al.}(2019)\citenamefont {Aurate}
  \emph {et~al.}}]{arute19}%
  \BibitemOpen
  \bibfield  {author} {\bibinfo {author} {\bibfnamefont {F.}~\bibnamefont
  {Aurate}} \emph {et~al.},\ }\bibfield  {title} {\bibinfo {title} {Quantum
  supremacy using a programmable superconducting processor},\ }\href
  {https://doi.org/10.1038/s41586-019-1666-5} {\bibfield  {journal} {\bibinfo
  {journal} {Nature}\ }\textbf {\bibinfo {volume} {574}},\ \bibinfo {pages}
  {505} (\bibinfo {year} {2019})}\BibitemShut {NoStop}%
\bibitem [{\citenamefont {Temme}\ \emph {et~al.}(2017)\citenamefont {Temme},
  \citenamefont {Bravyi},\ and\ \citenamefont {Gambetta}}]{temme17}%
  \BibitemOpen
  \bibfield  {author} {\bibinfo {author} {\bibfnamefont {K.}~\bibnamefont
  {Temme}}, \bibinfo {author} {\bibfnamefont {S.}~\bibnamefont {Bravyi}},\ and\
  \bibinfo {author} {\bibfnamefont {J.~M.}\ \bibnamefont {Gambetta}},\
  }\bibfield  {title} {\bibinfo {title} {Error mitigation for short-depth
  quantum circuits},\ }\href {https://doi.org/10.1103/PhysRevLett.119.180509}
  {\bibfield  {journal} {\bibinfo  {journal} {Phys. Rev. Lett.}\ }\textbf
  {\bibinfo {volume} {119}},\ \bibinfo {pages} {180509} (\bibinfo {year}
  {2017})}\BibitemShut {NoStop}%
\bibitem [{\citenamefont {van~den Berg}\ \emph {et~al.}(2023)\citenamefont
  {van~den Berg}, \citenamefont {Minev}, \citenamefont {Kandala},\ and\
  \citenamefont {Temme}}]{vandenberg23}%
  \BibitemOpen
  \bibfield  {author} {\bibinfo {author} {\bibfnamefont {E.}~\bibnamefont
  {van~den Berg}}, \bibinfo {author} {\bibfnamefont {Z.~K.}\ \bibnamefont
  {Minev}}, \bibinfo {author} {\bibfnamefont {A.}~\bibnamefont {Kandala}},\
  and\ \bibinfo {author} {\bibfnamefont {K.}~\bibnamefont {Temme}},\ }\bibfield
   {title} {\bibinfo {title} {Probabilistic error cancellation with sparse
  pauli--lindblad models on noisy quantum processors},\ }\bibfield  {journal}
  {\bibinfo  {journal} {Nature Physics}\ }\href
  {https://doi.org/10.1038/s41567-023-02042-2} {10.1038/s41567-023-02042-2}
  (\bibinfo {year} {2023})\BibitemShut {NoStop}%
\bibitem [{\citenamefont {Kim}\ \emph {et~al.}(2023)\citenamefont {Kim},
  \citenamefont {Eddins}, \citenamefont {Anand}, \citenamefont {Wei},
  \citenamefont {van~den Berg}, \citenamefont {Rosenblatt}, \citenamefont
  {Nayfeh}, \citenamefont {Wu}, \citenamefont {Zaletel}, \citenamefont
  {Temme},\ and\ \citenamefont {Kandala}}]{kim2023}%
  \BibitemOpen
  \bibfield  {author} {\bibinfo {author} {\bibfnamefont {Y.}~\bibnamefont
  {Kim}}, \bibinfo {author} {\bibfnamefont {A.}~\bibnamefont {Eddins}},
  \bibinfo {author} {\bibfnamefont {S.}~\bibnamefont {Anand}}, \bibinfo
  {author} {\bibfnamefont {K.~X.}\ \bibnamefont {Wei}}, \bibinfo {author}
  {\bibfnamefont {E.}~\bibnamefont {van~den Berg}}, \bibinfo {author}
  {\bibfnamefont {S.}~\bibnamefont {Rosenblatt}}, \bibinfo {author}
  {\bibfnamefont {H.}~\bibnamefont {Nayfeh}}, \bibinfo {author} {\bibfnamefont
  {Y.}~\bibnamefont {Wu}}, \bibinfo {author} {\bibfnamefont {M.}~\bibnamefont
  {Zaletel}}, \bibinfo {author} {\bibfnamefont {K.}~\bibnamefont {Temme}},\
  and\ \bibinfo {author} {\bibfnamefont {A.}~\bibnamefont {Kandala}},\
  }\bibfield  {title} {\bibinfo {title} {Evidence for the utility of quantum
  computing before fault tolerance},\ }\href
  {https://doi.org/10.1038/s41586-023-06096-3} {\bibfield  {journal} {\bibinfo
  {journal} {Nature}\ }\textbf {\bibinfo {volume} {618}},\ \bibinfo {pages}
  {500} (\bibinfo {year} {2023})}\BibitemShut {NoStop}%
\bibitem [{\citenamefont {Vidal}(2003)}]{Vidal2003}%
  \BibitemOpen
  \bibfield  {author} {\bibinfo {author} {\bibfnamefont {G.}~\bibnamefont
  {Vidal}},\ }\bibfield  {title} {\bibinfo {title} {Efficient {Classical}
  {Simulation} of {Slightly} {Entangled} {Quantum} {Computations}},\ }\href
  {https://doi.org/10.1103/PhysRevLett.91.147902} {\bibfield  {journal}
  {\bibinfo  {journal} {Physical Review Letters}\ }\textbf {\bibinfo {volume}
  {91}},\ \bibinfo {pages} {147902} (\bibinfo {year} {2003})}\BibitemShut
  {NoStop}%
\bibitem [{\citenamefont {Jiang}\ \emph {et~al.}(2008)\citenamefont {Jiang},
  \citenamefont {Weng},\ and\ \citenamefont {Xiang}}]{Jiang2008}%
  \BibitemOpen
  \bibfield  {author} {\bibinfo {author} {\bibfnamefont {H.~C.}\ \bibnamefont
  {Jiang}}, \bibinfo {author} {\bibfnamefont {Z.~Y.}\ \bibnamefont {Weng}},\
  and\ \bibinfo {author} {\bibfnamefont {T.}~\bibnamefont {Xiang}},\ }\bibfield
   {title} {\bibinfo {title} {Accurate {Determination} of {Tensor} {Network}
  {State} of {Quantum} {Lattice} {Models} in {Two} {Dimensions}},\ }\href
  {https://doi.org/10.1103/PhysRevLett.101.090603} {\bibfield  {journal}
  {\bibinfo  {journal} {Physical Review Letters}\ }\textbf {\bibinfo {volume}
  {101}},\ \bibinfo {pages} {090603} (\bibinfo {year} {2008})}\BibitemShut
  {NoStop}%
\bibitem [{\citenamefont {Sahu}\ and\ \citenamefont
  {Swingle}(2022)}]{TensorNetworkBeliefPropagation2}%
  \BibitemOpen
  \bibfield  {author} {\bibinfo {author} {\bibfnamefont {S.}~\bibnamefont
  {Sahu}}\ and\ \bibinfo {author} {\bibfnamefont {B.}~\bibnamefont {Swingle}},\
  }\href@noop {} {\bibinfo {title} {Efficient tensor network simulation of
  quantum many-body physics on sparse graphs}} (\bibinfo {year} {2022}),\
  \Eprint {https://arxiv.org/abs/2206.04701} {arXiv:2206.04701 [quant-ph]}
  \BibitemShut {NoStop}%
\bibitem [{\citenamefont {Tindall}\ and\ \citenamefont
  {Fishman}(2023)}]{TindallGauging2023}%
  \BibitemOpen
  \bibfield  {author} {\bibinfo {author} {\bibfnamefont {J.}~\bibnamefont
  {Tindall}}\ and\ \bibinfo {author} {\bibfnamefont {M.}~\bibnamefont
  {Fishman}},\ }\bibfield  {title} {\bibinfo {title} {Gauging tensor networks
  with belief propagation},\ }\href@noop {} {\bibfield  {journal} {\bibinfo
  {journal} {Preprint forthcoming}\ } (\bibinfo {year} {2023})}\BibitemShut
  {NoStop}%
\bibitem [{\citenamefont {Ran}\ \emph {et~al.}(2012)\citenamefont {Ran},
  \citenamefont {Li}, \citenamefont {Xi}, \citenamefont {Zhang},\ and\
  \citenamefont {Su}}]{Ran2012}%
  \BibitemOpen
  \bibfield  {author} {\bibinfo {author} {\bibfnamefont {S.-J.}\ \bibnamefont
  {Ran}}, \bibinfo {author} {\bibfnamefont {W.}~\bibnamefont {Li}}, \bibinfo
  {author} {\bibfnamefont {B.}~\bibnamefont {Xi}}, \bibinfo {author}
  {\bibfnamefont {Z.}~\bibnamefont {Zhang}},\ and\ \bibinfo {author}
  {\bibfnamefont {G.}~\bibnamefont {Su}},\ }\bibfield  {title} {\bibinfo
  {title} {Optimized decimation of tensor networks with super-orthogonalization
  for two-dimensional quantum lattice models},\ }\href
  {https://doi.org/10.1103/PhysRevB.86.134429} {\bibfield  {journal} {\bibinfo
  {journal} {Physical Review B}\ }\textbf {\bibinfo {volume} {86}},\ \bibinfo
  {pages} {134429} (\bibinfo {year} {2012})}\BibitemShut {NoStop}%
\bibitem [{\citenamefont {Phien}\ \emph {et~al.}(2015)\citenamefont {Phien},
  \citenamefont {McCulloch},\ and\ \citenamefont {Vidal}}]{Phien2015}%
  \BibitemOpen
  \bibfield  {author} {\bibinfo {author} {\bibfnamefont {H.~N.}\ \bibnamefont
  {Phien}}, \bibinfo {author} {\bibfnamefont {I.~P.}\ \bibnamefont
  {McCulloch}},\ and\ \bibinfo {author} {\bibfnamefont {G.}~\bibnamefont
  {Vidal}},\ }\bibfield  {title} {\bibinfo {title} {Fast convergence of
  imaginary time evolution tensor network algorithms by recycling the
  environment},\ }\href {https://doi.org/10.1103/PhysRevB.91.115137} {\bibfield
   {journal} {\bibinfo  {journal} {Physical Review B}\ }\textbf {\bibinfo
  {volume} {91}},\ \bibinfo {pages} {115137} (\bibinfo {year}
  {2015})}\BibitemShut {NoStop}%
\bibitem [{\citenamefont {Jahromi}\ and\ \citenamefont
  {Orús}(2019)}]{Jahromi2019}%
  \BibitemOpen
  \bibfield  {author} {\bibinfo {author} {\bibfnamefont {S.~S.}\ \bibnamefont
  {Jahromi}}\ and\ \bibinfo {author} {\bibfnamefont {R.}~\bibnamefont
  {Orús}},\ }\bibfield  {title} {\bibinfo {title} {Universal tensor-network
  algorithm for any infinite lattice},\ }\href
  {https://doi.org/10.1103/PhysRevB.99.195105} {\bibfield  {journal} {\bibinfo
  {journal} {Physical Review B}\ }\textbf {\bibinfo {volume} {99}},\ \bibinfo
  {pages} {195105} (\bibinfo {year} {2019})}\BibitemShut {NoStop}%
\bibitem [{\citenamefont {Alkabetz}\ and\ \citenamefont
  {Arad}(2021)}]{SimpleUpdateGauging3}%
  \BibitemOpen
  \bibfield  {author} {\bibinfo {author} {\bibfnamefont {R.}~\bibnamefont
  {Alkabetz}}\ and\ \bibinfo {author} {\bibfnamefont {I.}~\bibnamefont
  {Arad}},\ }\bibfield  {title} {\bibinfo {title} {Tensor networks contraction
  and the belief propagation algorithm},\ }\href
  {https://doi.org/10.1103/PhysRevResearch.3.023073} {\bibfield  {journal}
  {\bibinfo  {journal} {Phys. Rev. Res.}\ }\textbf {\bibinfo {volume} {3}},\
  \bibinfo {pages} {023073} (\bibinfo {year} {2021})}\BibitemShut {NoStop}%
\bibitem [{\citenamefont {Verstraete}\ and\ \citenamefont
  {Cirac}(2004)}]{Verstraete2004}%
  \BibitemOpen
  \bibfield  {author} {\bibinfo {author} {\bibfnamefont {F.}~\bibnamefont
  {Verstraete}}\ and\ \bibinfo {author} {\bibfnamefont {J.~I.}\ \bibnamefont
  {Cirac}},\ }\href {https://doi.org/10.48550/arXiv.cond-mat/0407066} {\emph
  {\bibinfo {title} {Renormalization algorithms for {Quantum}-{Many} {Body}
  {Systems} in two and higher dimensions}}},\ \bibinfo {type} {Tech. Rep.}\
  (\bibinfo {year} {2004})\ \bibinfo {note} {arXiv:cond-mat/0407066 type:
  article}\BibitemShut {NoStop}%
\bibitem [{\citenamefont {Chubb}(2021)}]{chubb2021general}%
  \BibitemOpen
  \bibfield  {author} {\bibinfo {author} {\bibfnamefont {C.~T.}\ \bibnamefont
  {Chubb}},\ }\href@noop {} {\bibinfo {title} {General tensor network decoding
  of 2d pauli codes}} (\bibinfo {year} {2021}),\ \Eprint
  {https://arxiv.org/abs/2101.04125} {arXiv:2101.04125 [quant-ph]} \BibitemShut
  {NoStop}%
\bibitem [{\citenamefont {Zaletel}\ and\ \citenamefont
  {Pollmann}(2020)}]{Zaletel2020}%
  \BibitemOpen
  \bibfield  {author} {\bibinfo {author} {\bibfnamefont {M.~P.}\ \bibnamefont
  {Zaletel}}\ and\ \bibinfo {author} {\bibfnamefont {F.}~\bibnamefont
  {Pollmann}},\ }\bibfield  {title} {\bibinfo {title} {Isometric {Tensor}
  {Network} {States} in {Two} {Dimensions}},\ }\href
  {https://doi.org/10.1103/PhysRevLett.124.037201} {\bibfield  {journal}
  {\bibinfo  {journal} {Physical Review Letters}\ }\textbf {\bibinfo {volume}
  {124}},\ \bibinfo {pages} {037201} (\bibinfo {year} {2020})}\BibitemShut
  {NoStop}%
\bibitem [{\citenamefont {Soejima}\ \emph {et~al.}(2020)\citenamefont
  {Soejima}, \citenamefont {Siva}, \citenamefont {Bultinck}, \citenamefont
  {Chatterjee}, \citenamefont {Pollmann},\ and\ \citenamefont
  {Zaletel}}]{Soejima2020}%
  \BibitemOpen
  \bibfield  {author} {\bibinfo {author} {\bibfnamefont {T.}~\bibnamefont
  {Soejima}}, \bibinfo {author} {\bibfnamefont {K.}~\bibnamefont {Siva}},
  \bibinfo {author} {\bibfnamefont {N.}~\bibnamefont {Bultinck}}, \bibinfo
  {author} {\bibfnamefont {S.}~\bibnamefont {Chatterjee}}, \bibinfo {author}
  {\bibfnamefont {F.}~\bibnamefont {Pollmann}},\ and\ \bibinfo {author}
  {\bibfnamefont {M.~P.}\ \bibnamefont {Zaletel}},\ }\bibfield  {title}
  {\bibinfo {title} {Isometric tensor network representation of string-net
  liquids},\ }\href {https://doi.org/10.1103/PhysRevB.101.085117} {\bibfield
  {journal} {\bibinfo  {journal} {Physical Review B}\ }\textbf {\bibinfo
  {volume} {101}},\ \bibinfo {pages} {085117} (\bibinfo {year}
  {2020})}\BibitemShut {NoStop}%
\bibitem [{\citenamefont {Murg}\ \emph {et~al.}(2007)\citenamefont {Murg},
  \citenamefont {Verstraete},\ and\ \citenamefont {Cirac}}]{Murg2007}%
  \BibitemOpen
  \bibfield  {author} {\bibinfo {author} {\bibfnamefont {V.}~\bibnamefont
  {Murg}}, \bibinfo {author} {\bibfnamefont {F.}~\bibnamefont {Verstraete}},\
  and\ \bibinfo {author} {\bibfnamefont {J.~I.}\ \bibnamefont {Cirac}},\
  }\bibfield  {title} {\bibinfo {title} {Variational study of hard-core bosons
  in a two-dimensional optical lattice using projected entangled pair states},\
  }\href {https://doi.org/10.1103/PhysRevA.75.033605} {\bibfield  {journal}
  {\bibinfo  {journal} {Physical Review A}\ }\textbf {\bibinfo {volume} {75}},\
  \bibinfo {pages} {033605} (\bibinfo {year} {2007})}\BibitemShut {NoStop}%
\bibitem [{\citenamefont {Tindall}\ \emph {et~al.}(2022)\citenamefont
  {Tindall}, \citenamefont {Searle}, \citenamefont {Alhajri},\ and\
  \citenamefont {Jaksch}}]{tindall2022quantum}%
  \BibitemOpen
  \bibfield  {author} {\bibinfo {author} {\bibfnamefont {J.}~\bibnamefont
  {Tindall}}, \bibinfo {author} {\bibfnamefont {A.}~\bibnamefont {Searle}},
  \bibinfo {author} {\bibfnamefont {A.}~\bibnamefont {Alhajri}},\ and\ \bibinfo
  {author} {\bibfnamefont {D.}~\bibnamefont {Jaksch}},\ }\bibfield  {title}
  {\bibinfo {title} {Quantum physics in connected worlds},\ }\href@noop {}
  {\bibfield  {journal} {\bibinfo  {journal} {Nature Communications}\ }\textbf
  {\bibinfo {volume} {13}},\ \bibinfo {pages} {7445} (\bibinfo {year}
  {2022})}\BibitemShut {NoStop}%
\bibitem [{\citenamefont {Searle}\ and\ \citenamefont
  {Tindall}(2023)}]{searle2023exact}%
  \BibitemOpen
  \bibfield  {author} {\bibinfo {author} {\bibfnamefont {A.}~\bibnamefont
  {Searle}}\ and\ \bibinfo {author} {\bibfnamefont {J.}~\bibnamefont
  {Tindall}},\ }\href@noop {} {\bibinfo {title} {An exact continuous theory for
  spin systems on graphons}} (\bibinfo {year} {2023}),\ \Eprint
  {https://arxiv.org/abs/2303.00731} {arXiv:2303.00731 [cond-mat.stat-mech]}
  \BibitemShut {NoStop}%
\bibitem [{ITe(2023)}]{ITensorNetworks}%
  \BibitemOpen
  \href@noop {} {\bibinfo {title} {{ITensorNetworks.jl}}},\ \bibinfo
  {howpublished} {\url{https://github.com/mtfishman/ITensorNetworks.jl}}
  (\bibinfo {year} {2023})\BibitemShut {NoStop}%
\bibitem [{\citenamefont {Wang}\ \emph {et~al.}(2011)\citenamefont {Wang},
  \citenamefont {Pižorn},\ and\ \citenamefont {Verstraete}}]{Wang2011}%
  \BibitemOpen
  \bibfield  {author} {\bibinfo {author} {\bibfnamefont {L.}~\bibnamefont
  {Wang}}, \bibinfo {author} {\bibfnamefont {I.}~\bibnamefont {Pižorn}},\ and\
  \bibinfo {author} {\bibfnamefont {F.}~\bibnamefont {Verstraete}},\ }\bibfield
   {title} {\bibinfo {title} {Monte {Carlo} simulation with tensor network
  states},\ }\href {https://doi.org/10.1103/PhysRevB.83.134421} {\bibfield
  {journal} {\bibinfo  {journal} {Physical Review B}\ }\textbf {\bibinfo
  {volume} {83}},\ \bibinfo {pages} {134421} (\bibinfo {year}
  {2011})}\BibitemShut {NoStop}%
\bibitem [{\citenamefont {Li}\ \emph {et~al.}(2012)\citenamefont {Li},
  \citenamefont {von Delft},\ and\ \citenamefont {Xiang}}]{Li2012}%
  \BibitemOpen
  \bibfield  {author} {\bibinfo {author} {\bibfnamefont {W.}~\bibnamefont
  {Li}}, \bibinfo {author} {\bibfnamefont {J.}~\bibnamefont {von Delft}},\ and\
  \bibinfo {author} {\bibfnamefont {T.}~\bibnamefont {Xiang}},\ }\bibfield
  {title} {\bibinfo {title} {Efficient simulation of infinite tree tensor
  network states on the {Bethe} lattice},\ }\href
  {https://doi.org/10.1103/PhysRevB.86.195137} {\bibfield  {journal} {\bibinfo
  {journal} {Physical Review B}\ }\textbf {\bibinfo {volume} {86}},\ \bibinfo
  {pages} {195137} (\bibinfo {year} {2012})}\BibitemShut {NoStop}%
\bibitem [{\citenamefont {Pearl}(1982)}]{BeliefPropagation1}%
  \BibitemOpen
  \bibfield  {author} {\bibinfo {author} {\bibfnamefont {J.}~\bibnamefont
  {Pearl}},\ }\bibfield  {title} {\bibinfo {title} {Reverend bayes on inference
  engines: A distributed hierarchical approach},\ }in\ \href@noop {} {\emph
  {\bibinfo {booktitle} {Proceedings of the Second AAAI Conference on
  Artificial Intelligence}}},\ \bibinfo {series and number} {AAAI'82}\
  (\bibinfo  {publisher} {AAAI Press},\ \bibinfo {year} {1982})\ p.\ \bibinfo
  {pages} {133–136}\BibitemShut {NoStop}%
\bibitem [{\citenamefont {Leifer}\ and\ \citenamefont
  {Poulin}(2008)}]{TensorNetworkBeliefPropagation1}%
  \BibitemOpen
  \bibfield  {author} {\bibinfo {author} {\bibfnamefont {M.}~\bibnamefont
  {Leifer}}\ and\ \bibinfo {author} {\bibfnamefont {D.}~\bibnamefont
  {Poulin}},\ }\bibfield  {title} {\bibinfo {title} {Quantum graphical models
  and belief propagation},\ }\href
  {https://doi.org/https://doi.org/10.1016/j.aop.2007.10.001} {\bibfield
  {journal} {\bibinfo  {journal} {Annals of Physics}\ }\textbf {\bibinfo
  {volume} {323}},\ \bibinfo {pages} {1899} (\bibinfo {year}
  {2008})}\BibitemShut {NoStop}%
\bibitem [{\citenamefont {Guo}\ \emph {et~al.}(2023)\citenamefont {Guo},
  \citenamefont {Poletti},\ and\ \citenamefont
  {Arad}}]{TensorNetworkBeliefPropagation3}%
  \BibitemOpen
  \bibfield  {author} {\bibinfo {author} {\bibfnamefont {C.}~\bibnamefont
  {Guo}}, \bibinfo {author} {\bibfnamefont {D.}~\bibnamefont {Poletti}},\ and\
  \bibinfo {author} {\bibfnamefont {I.}~\bibnamefont {Arad}},\ }\href@noop {}
  {\bibinfo {title} {Block belief propagation algorithm for 2d tensor
  networks}} (\bibinfo {year} {2023}),\ \Eprint
  {https://arxiv.org/abs/2301.05844} {arXiv:2301.05844 [quant-ph]} \BibitemShut
  {NoStop}%
\bibitem [{\citenamefont {Bezanson}\ \emph {et~al.}(2017)\citenamefont
  {Bezanson}, \citenamefont {Edelman}, \citenamefont {Karpinski},\ and\
  \citenamefont {Shah}}]{bezanson2017julia}%
  \BibitemOpen
  \bibfield  {author} {\bibinfo {author} {\bibfnamefont {J.}~\bibnamefont
  {Bezanson}}, \bibinfo {author} {\bibfnamefont {A.}~\bibnamefont {Edelman}},
  \bibinfo {author} {\bibfnamefont {S.}~\bibnamefont {Karpinski}},\ and\
  \bibinfo {author} {\bibfnamefont {V.~B.}\ \bibnamefont {Shah}},\ }\bibfield
  {title} {\bibinfo {title} {Julia: A fresh approach to numerical computing},\
  }\href {https://doi.org/https://doi.org/10.1137/141000671} {\bibfield
  {journal} {\bibinfo  {journal} {SIAM review}\ }\textbf {\bibinfo {volume}
  {59}},\ \bibinfo {pages} {65} (\bibinfo {year} {2017})}\BibitemShut {NoStop}%
\bibitem [{\citenamefont {Fishman}\ \emph {et~al.}(2022)\citenamefont
  {Fishman}, \citenamefont {White},\ and\ \citenamefont
  {Stoudenmire}}]{itensor-r0.3}%
  \BibitemOpen
  \bibfield  {author} {\bibinfo {author} {\bibfnamefont {M.}~\bibnamefont
  {Fishman}}, \bibinfo {author} {\bibfnamefont {S.~R.}\ \bibnamefont {White}},\
  and\ \bibinfo {author} {\bibfnamefont {E.~M.}\ \bibnamefont {Stoudenmire}},\
  }\bibfield  {title} {\bibinfo {title} {{Codebase release 0.3 for ITensor}},\
  }\href {https://doi.org/10.21468/SciPostPhysCodeb.4-r0.3} {\bibfield
  {journal} {\bibinfo  {journal} {SciPost Phys. Codebases}\ ,\ \bibinfo {pages}
  {4}} (\bibinfo {year} {2022})}\BibitemShut {NoStop}%
\bibitem [{Gra(2023)}]{GraphTikz}%
  \BibitemOpen
  \href@noop {} {\bibinfo {title} {{GraphTikz.jl}}},\ \bibinfo {howpublished}
  {\url{https://github.com/mtfishman/GraphTikZ.jl}} (\bibinfo {year}
  {2023})\BibitemShut {NoStop}%
\end{thebibliography}%

\end{document}
